\chapter{Родитель индуцированной ньютоновской постоянной}
\label{ch:v10-induced-newton-parent}

\section{Эйнштейновский коэффициент и ньютоновская площадь}

В безразмерном действии коэффициент эйнштейновского члена удобно связывать
непосредственно с геометрической ньютоновской площадью
\begin{equation}
 g_N:=\frac{\hbar G}{c^3},
 \qquad 16\pi g_N B=1.
 \label{eq:v10-newton-einstein-normalization}
\end{equation}
Следовательно,
\begin{equation}
 g_N=\frac1{16\pi B},
 \qquad
 G=\frac{c^3}{16\pi\hbar B}.
 \label{eq:v10-newton-from-B}
\end{equation}

Предыдущий кривизненный сектор допускал индуцированную форму
\begin{equation}
 A=\alpha m^2,
 \qquad B=\beta m,
 \qquad
 q:=\ell_{\mathrm{cell}}^2=\frac{\beta}{2\alpha m}.
 \label{eq:v10-newton-induced-data}
\end{equation}
Поэтому условно
\begin{equation}
 g_N=\frac1{16\pi\beta m},
 \qquad
 G=\frac{c^3}{16\pi\hbar\beta m}.
 \label{eq:v10-newton-induced-constant}
\end{equation}

\section{Слепое отношение ньютоновской и клеточной площадей}

Хотя абсолютная величина $m$ остаётся свободной, она сокращается в отношении
\begin{equation}
 \boxed{\frac{g_N}{\ell_{\mathrm{cell}}^2}
 =\frac{g_N}{q}
 =\frac{\alpha}{8\pi\beta^2}}.
 \label{eq:v10-newton-cell-blind-ratio}
\end{equation}
Это новое безразмерное следствие: один и тот же индуцированный масштаб
контролирует ньютоновскую и клеточную площади.

Три нормировки можно собрать в положительный цепочечный родитель:
\begin{equation}
 u_E=16\pi g_NB,
 \qquad u_m=\frac{B}{\beta m},
 \qquad u_q=\frac{2\alpha qm}{\beta},
\end{equation}
\begin{equation}
 \mathcal P_N
 =\frac12(u_E-1)^2+\frac12(u_m-u_E)^2+\frac12(u_q-u_m)^2.
 \label{eq:v10-newton-common-parent}
\end{equation}
Его единственный нуль равен $u_E=u_m=u_q=1$. Гессиан имеет вид
\begin{equation}
 \begin{pmatrix}2&-1&0\\-1&2&-1\\0&-1&1\end{pmatrix},
 \qquad \rank=3,
 \qquad \det=1,
\end{equation}
а ведущие главные миноры равны $(2,3,1)$.

\section{Оставшаяся масштабная мода}

В логарифмических переменных $(g_N,B,m,q)$ три условия имеют матрицу
\begin{equation}
 C_N=
 \begin{pmatrix}
 1&1&0&0\\
 0&1&-1&0\\
 0&0&1&1
 \end{pmatrix}.
 \label{eq:v10-newton-scale-map}
\end{equation}
Она имеет ранг/ядро $3/1$ и уничтожает $(1,-1,-1,1)^T$. Поэтому
\begin{equation}
 (g_N,B,m,q)\mapsto(s^2g_N,B/s^2,m/s^2,s^2q)
\end{equation}
сохраняет весь условный родитель. Эйнштейновская нормировка переводит
свободу коэффициента $B$ в свободу $G$, но не уничтожает её.

\begin{theorem}[Относительное происхождение ньютоновской площади]
При условных коэффициентах $A=\alpha m^2$ и $B=\beta m$ один положительный
родитель совместно выбирает эйнштейновскую нормировку, коэффициент и размер
ячейки. Отношение $g_N/\ell_{\mathrm{cell}}^2=\alpha/(8\pi\beta^2)$ не зависит
от размерного семени $m$.
\end{theorem}

\begin{nogocriterion}[Индуцированный коэффициент не является абсолютным якорем]
Формула $G=c^3/(16\pi\hbar\beta m)$ не выводит ньютоновскую постоянную,
пока общий родитель не выбирает абсолютную величину $m$ и коэффициент
$\beta$. Общая масштабная мода одновременно растягивает клеточную и
ньютоновскую площади.
\end{nogocriterion}

\begin{protocol}[Статус ньютоновской постоянной]\begin{enumerate}
\item эйнштейновская нормировка: \textbf{$1/1$};
\item положительный относительный родитель: \textbf{построен};
\item относительное происхождение: \textbf{$6/6$};
\item слепое отношение площадей: \textbf{$1/1$};
\item происхождение $m$ и коэффициентов: \textbf{$0/2$};
\item абсолютная ньютоновская постоянная: \textbf{$0/1$};
\item следующий гейт: \textbf{аудит источников размерного семени $m$}.
\end{enumerate}\end{protocol}

Машинный сертификат сохранён в
\path{s2t_v10_cell_birth_four_volume_induced_newton_constant_parent_origin_gate_results.json}.